
\pagebreak

\chapter{Technische Grundlagen}
\abschnittsautor{Arslan Yunus}

\section{Signaturbasierte Erkennung und maschinelles Lernen}

Dieser Abschnitt erklärt die Grundlagen der ersten Erkennungsphase des
Systems: wie Dateien anhand von Hashes, Mustern und maschinellem Lernen
als gut- oder bösartig eingestuft werden.

\subsection{Das PE-Dateiformat}

Ausführbare Dateien unter Windows liegen im \ac{PE}-Format vor
(Portable Executable). Eine solche Datei besteht aus einem Header mit
Metadaten (z.\,B. Zielarchitektur, Einsprungpunkt) sowie mehreren
Sections, etwa \texttt{.text} für den ausführbaren Code oder
\texttt{.rsrc} für Ressourcen. Der Header enthält außerdem eine
Import-Tabelle, die auflistet, welche externen Funktionen die Datei beim
Start nachlädt~\cite{pe-format}. Diese Struktur ist wichtig, weil genau
diese Informationen später als Merkmale zur Erkennung von Schadsoftware
genutzt werden.

\subsection{Kryptografische Hashfunktionen}

Eine Hashfunktion wie \ac{SHA}-256 wandelt eine beliebig lange Datei in
einen festen, kurzen Wert um. Schon eine winzige Änderung an der Datei
(ein einziges verändertes Byte) erzeugt dabei einen komplett anderen
Hashwert~\cite{sha256-spec}. Dadurch lassen sich bekannte
Malware-Dateien eindeutig identifizieren, indem man ihren Hash mit einer
Datenbank bekannter Schadsoftware vergleicht. Das Problem: Ändert ein
Angreifer auch nur eine Kleinigkeit am Code, ändert sich der Hash
komplett und die Erkennung schlägt fehl.

\subsection{Fuzzy Hashing mit TLSH}

\ac{TLSH} (Trend Micro Locality Sensitive Hash) löst genau dieses
Problem: Im Gegensatz zu \ac{SHA}-256 erzeugt \ac{TLSH} für ähnliche
Dateien auch ähnliche Hashwerte. Man kann also zwei Hashes vergleichen
und eine Zahl bekommen, die angibt, wie ähnlich sich die Originaldateien
sind~\cite{tlsh-paper}. Dadurch erkennt das System auch leicht
abgewandelte Varianten einer bekannten Malware-Familie, selbst wenn der
klassische Hash längst nicht mehr passt.

\subsection{Regelbasierte Mustererkennung mit YARA}

\ac{YARA} erlaubt es, eigene Erkennungsregeln zu schreiben: Eine Regel
definiert eine Liste von verdächtigen Text- oder Byte-Mustern (z.\,B.
bestimmte Zeichenketten oder Code-Fragmente) sowie eine Bedingung, wie
viele davon in einer Datei vorkommen müssen, damit sie als schädlich
gilt~\cite{yara-docs}. So lassen sich ganze Malware-Familien an
typischen Merkmalen erkennen, auch wenn sich einzelne Dateien im Detail
unterscheiden.

\subsection{Whitelisting vertrauenswürdiger Software}

Neben der Erkennung bösartiger Dateien ist das gezielte Freigeben
bekannter, vertrauenswürdiger Software ein ebenso wichtiger Baustein der
signaturbasierten Phase. Beim Whitelisting wird der Hash (siehe
vorheriger Abschnitt) einer Datei gegen eine Liste bekannter, geprüfter
Software abgeglichen -- stammt eine Datei nachweislich von einem
vertrauenswürdigen Hersteller, wird sie von der weiteren, rechenintensiveren
Analyse (Verhaltensüberwachung, Klassifikation) ausgenommen. Das reduziert
sowohl die Systemlast als auch die Zahl potenzieller Fehlalarme, da
bekanntermaßen harmlose Systemdateien und verbreitete Anwendungen nicht
bei jedem Start erneut vollständig analysiert werden müssen. Eine robuste
Whitelist ist damit eine notwendige Ergänzung zur reinen
Blacklist-Erkennung, denn ohne sie würde ein rein verhaltens- oder
KI-basiertes System deutlich mehr falsch-positive Meldungen erzeugen.

\subsection{Multi-Engine-Abgleich mit VirusTotal}

Ergänzend zur eigenen Hash- und Mustererkennung nutzt das System die
\ac{API} von VirusTotal, einem Online-Dienst, der eine eingereichte
Datei oder deren Hash gegen eine Vielzahl unabhängiger Antiviren-Engines
gleichzeitig prüft und ein aggregiertes Ergebnis zurückliefert, wie viele
der Engines die Datei als bösartig einstufen~\cite{virustotal-docs}. Der
entscheidende Vorteil gegenüber einer einzelnen, lokal betriebenen
Engine liegt in der Konsens-Bildung: Stufen mehrere unabhängige Anbieter
eine Datei übereinstimmend als schädlich ein, erhöht das die
Verlässlichkeit der Entscheidung deutlich gegenüber einer Einzelmeinung.
Da die kostenlose Public API auf wenige Anfragen pro Minute limitiert
ist, wird sie im System primär als ergänzende Bestätigungsinstanz
eingesetzt und nicht als alleinige Erkennungsgrundlage -- die
eigentliche Klassifikation erfolgt weiterhin lokal über die im Folgenden
beschriebenen Modelle, um die Datenhoheit zu wahren und nicht bei jeder
Datei auf eine externe Cloud-Anfrage angewiesen zu sein.

\subsection{Gradient Boosting: XGBoost und LightGBM}

Für die eigentliche Klassifikation (gutartig vs. bösartig) kommt
maschinelles Lernen zum Einsatz, konkret \ac{XGBoost} und
\ac{LightGBM}. Beide funktionieren nach dem Prinzip des Gradient
Boosting: Es wird eine Reihe einfacher Entscheidungsbäume nacheinander
trainiert, wobei jeder neue Baum gezielt die Fehler der vorherigen Bäume
ausbessert~\cite{xgboost-paper}. \ac{LightGBM} ist dabei eine schnellere
Variante desselben Prinzips, optimiert für sehr große
Datenmengen~\cite{lightgbm-paper}. Beide Modelle eignen sich gut für
tabellarische Daten -- also genau die Art von Merkmalen, die man aus
einer \ac{PE}-Datei extrahieren kann.

\subsection{Merkmalsextraktion aus PE-Dateien}

Damit ein Modell wie \ac{XGBoost} überhaupt etwas lernen kann, muss jede
Datei zuerst in Zahlen umgewandelt werden. Typische Merkmale sind die
Entropie einzelner Sections (ein Maß dafür, wie "zufällig" die Daten
aussehen -- hohe Werte deuten oft auf verschlüsselten oder gepackten
Code hin), die Liste der importierten Funktionen sowie
Header-Informationen~\cite{ember-paper}. Diese Merkmale bilden gemeinsam
den Eingabevektor für das Modell und ersetzen die starre
Signaturprüfung durch eine lernbasierte, verallgemeinerbare Erkennung.

\pagebreak 
\abschnittsautor{Dervishi Drilon}
\section{Verhaltensbasierte Überwachung und Sequenzanalyse}

Dieser Abschnitt erklärt die Grundlagen der zweiten Erkennungsphase des
Systems: wie das Verhalten eines Programms zur Laufzeit beobachtet und
auf verdächtige Muster hin analysiert wird -- unabhängig davon, wie die
Datei selbst aussieht.

\subsection{Warum Verhaltensanalyse}

Signaturbasierte Methoden (siehe vorheriger Abschnitt) scheitern, sobald
Angreifer ihren Code gezielt verändern, um Hashes und Muster zu umgehen.
Schadsoftware muss jedoch, egal wie sie verpackt ist, zur Laufzeit
bestimmte Aktionen ausführen -- Dateien verändern, Prozesse starten,
Netzwerkverbindungen aufbauen. Die verhaltensbasierte Überwachung setzt
genau hier an: Statt die Datei selbst zu analysieren, wird beobachtet,
was ein Programm tut, nachdem es gestartet wurde.

\subsection{Event Tracing for Windows (ETW)}

\ac{ETW} ist ein in Windows eingebautes Protokollierungssystem, das
Systemereignisse wie Dateizugriffe, Prozessstarts oder
Registry-Änderungen in Echtzeit erfasst, ohne dass eine zusätzliche
Software tief ins Betriebssystem eingreifen muss~\cite{etw-docs}.
Programme (sogenannte Provider) senden Ereignisse an den ETW-Kernel, von
wo aus sie durch Consumer-Anwendungen ausgelesen werden können. Für das
System bedeutet das: ein performanter, bereits im Betriebssystem
vorhandener Kanal, über den sich verdächtiges Verhalten beobachten
lässt, ohne das System spürbar zu verlangsamen.

\subsection{Sysmon}

\ac{Sysmon} (System Monitor) ist ein kostenloses
Microsoft-Sysinternals-Tool, das auf \ac{ETW} aufbaut und detaillierte,
sicherheitsrelevante Ereignisse protokolliert -- etwa Prozesserstellung
mit vollständiger Kommandozeile, Netzwerkverbindungen oder Änderungen an
Autostart-Einträgen~\cite{sysmon-docs}. Über eine XML-Konfigurationsdatei
lässt sich genau festlegen, welche Ereignistypen protokolliert werden
sollen. Sysmon liefert damit die rohen, strukturierten Verhaltensdaten,
die im System als Grundlage für die spätere Sequenzanalyse dienen.

\subsection{API-Hooking mit Frida}

Während \ac{ETW} und \ac{Sysmon} Ereignisse auf Betriebssystemebene
protokollieren, erlaubt Frida einen noch tieferen Einblick: Es klinkt
sich direkt in den Prozessspeicher eines laufenden Programms ein und
kann einzelne Funktionsaufrufe (\ac{API}-Calls) abfangen, bevor oder
nachdem sie ausgeführt werden~\cite{frida-docs}. Dadurch lässt sich
beispielsweise genau mitverfolgen, mit welchen Parametern eine
verdächtige Funktion wie \texttt{VirtualAllocEx} aufgerufen wurde --
Informationen, die über \ac{ETW} allein nicht sichtbar wären. Frida
wird im System eingesetzt, um besonders relevante \ac{API}-Aufrufe
(z.\,B. Prozess- und Speicherfunktionen) mit zusätzlichem Kontext
anzureichern.

\subsection{CAPE Sandbox}

Da verdächtige Samples nicht auf einem produktiven System ausgeführt
werden dürfen, kommt \ac{CAPE} (Config And Payload Extraction) als
automatisierte Sandbox-Umgebung zum Einsatz. \ac{CAPE} führt eine Datei
in einer isolierten virtuellen Maschine aus, zeichnet dabei sämtliche
Systemaufrufe, Netzwerkaktivität und Prozessverhalten auf und kann
zusätzlich automatisch Konfigurationsdaten aus bekannten
Malware-Familien extrahieren~\cite{cape-sandbox}. Im System liefert
\ac{CAPE} damit vollständige, sichere Verhaltensprotokolle, die als
Trainingsgrundlage für das Sequenzmodell dienen.

\subsection{Long Short-Term Memory (LSTM)}

Die gesammelten Verhaltensdaten (Reihenfolge von \ac{API}-Aufrufen,
Prozess- und Dateiereignissen) bilden eine zeitliche Abfolge -- eine
Sequenz. Um solche Sequenzen zu klassifizieren, wird ein \ac{LSTM}
eingesetzt, eine spezielle Art von rekurrentem neuronalen Netz. Im
Unterschied zu einfachen neuronalen Netzen kann sich ein \ac{LSTM} über
viele Zeitschritte hinweg relevante Informationen "merken" und
irrelevante vergessen, gesteuert über sogenannte Gates (Input-, Forget-
und Output-Gate)~\cite{lstm-paper}. Dadurch eignet sich das Modell
besonders gut, um langfristige Verhaltensmuster zu erkennen -- etwa
wenn ein Programm erst nach mehreren Minuten unauffälligen Verhaltens
beginnt, verdächtige Aktionen durchzuführen (zeitverzögerte
Payload-Ausführung), was rein momentbasierte Methoden übersehen würden.

\subsection{Von Rohdaten zur Sequenz}

Bevor die \ac{LSTM}-Architektur trainiert werden kann, müssen die
rohen Ereignisse aus \ac{ETW}, \ac{Sysmon} und Frida in eine
einheitliche, zeitlich geordnete Sequenz überführt werden --
typischerweise als Folge von kategorisierten Ereignistypen (z.\,B.
ProcessCreate, FileWrite, NetworkConnect), die anschließend in
numerische Vektoren (Embeddings) umgewandelt werden, damit das Modell
damit rechnen kann. Diese Aufbereitung entspricht im Verhaltensbereich
dem, was die Merkmalsextraktion (Abschnitt zuvor) für statische Dateien
leistet: Sie macht rohe, heterogene Daten für ein lernbasiertes Modell
nutzbar.